% 本文件正文由作者撰写。
% !TeX root = main.tex
\section{结果}

\subsection{解码性能}
为评估单词级神经信号所包含的词汇可解码信息，我们首先比较了仅依赖当前词神经活动的 CNN 模型与进一步建模连续词神经上下文的 CNN+Transformer 模型。结果表明，在 ChineseEEG2、SMN4Lang 和 Pallier2025 三个数据集上，CNN-only 均能够从单词对应的神经信号中恢复高于机会水平的词汇信息，说明局部神经活动本身已经具有稳定的词汇可辨识性。

在此基础上，引入连续词之间的神经上下文后，CNN+Transformer 在三个数据集上均表现出一致的性能提升。该趋势同时出现在中文 EEG、中文 MEG 以及法语 MEG 数据中，表明上下文建模带来的增益并不局限于特定语言或神经记录模态。总体而言，单词级神经表征提供了解码的基本信息来源，而对连续神经活动之间关系的进一步建模能够显著增强词汇恢复能力。具体结果见图~\ref{fig:decoding_performance}。

\begin{figure}[htbp]
    \centering
    \includegraphics[width=\textwidth]{fig2_decoding_performance.pdf}
    \caption{CNN-only 与 CNN+Transformer 在三个数据集上的词汇解码性能。}
    \label{fig:decoding_performance}
\end{figure}



\subsection{审计}

为检验上下文模型的性能提升是否可能主要来自序列结构、位置或局部上下文先验，而非目标词对应的神经证据，我们进一步进行了严格对照实验，包括时间错位、同记录供体替换以及结构保留但神经特征置零等条件。结果表明，在三个数据集上，当全部神经特征被移除、仅保留分组、顺序与位置结构时，模型性能均下降至接近机会水平（图~\ref{fig:audit}）。这说明，仅依赖序列长度、位置和分组结构不足以解释上下文模型在正常条件下获得的高解码性能。

进一步地，当仅破坏目标词神经表征的时间对应关系，或将其替换为同一记录中其他词的神经表征时，解码性能均出现明显下降，而上下文结构本身保持不变。相比 CNN-only，上下文模型在这些扰动条件下表现出更大的性能损失，说明其较高的正常解码性能并非通过绕开目标神经信号获得，而是更强地依赖于正确对齐的神经表征。总体而言，这些结果支持这样一种解释：上下文建模主要增强了模型对与目标词一致的神经证据的利用，而不是简单地以序列结构或上下文先验替代神经信息。需要强调的是，这些对照反映的是不同条件下的描述性性能差异，并不构成对“纯神经信息”成分的可加性分解。

\begin{figure}[htbp]
    \centering
    \includegraphics[width=\textwidth]{fig3_audit.pdf}
    \caption{不同严格对照条件下的词汇解码性能。}
    \label{fig:audit}
\end{figure}
